\documentclass[final]{beamer}
\usepackage[size=custom,width=91.44,height=60.96,scale=1.04]{beamerposter}
\usepackage{amsmath,amssymb}
\usepackage{graphicx}
\usepackage{booktabs}
\usepackage{tikz}
\usetikzlibrary{arrows.meta,positioning,shapes.geometric}
\usepackage{ragged2e}
\usepackage{microtype}

% ---------- Theme ----------
\definecolor{Magenta}{HTML}{C2185B}
\definecolor{Orange}{HTML}{F57C00}
\definecolor{DeepPurple}{HTML}{5E2A84}
\definecolor{Green}{HTML}{2E7D32}
\definecolor{Red}{HTML}{C62828}
\definecolor{LightPink}{HTML}{FFF2F7}
\definecolor{LightOrange}{HTML}{FFF7EE}
\definecolor{LightGray}{HTML}{F5F5F5}
\definecolor{DarkText}{HTML}{222222}

\setbeamercolor{background canvas}{bg=white}
\setbeamercolor{normal text}{fg=DarkText,bg=white}
\setbeamercolor{block title}{fg=white,bg=Magenta}
\setbeamercolor{block body}{fg=DarkText,bg=white}
\setbeamercolor{block title alerted}{fg=white,bg=Orange}
\setbeamercolor{block body alerted}{fg=DarkText,bg=white}
\setbeamerfont{block title}{series=\bfseries,size=\large}
\setbeamerfont{block body}{size=\normalsize}
\setbeamertemplate{navigation symbols}{}
\setbeamertemplate{blocks}[rounded][shadow=false]

\newcommand{\tick}{\textcolor{Green}{\bfseries\checkmark}}
\newcommand{\cross}{\textcolor{Red}{\bfseries$\times$}}
\newcommand{\tightitem}{\setlength{\itemsep}{0.18em}\setlength{\parskip}{0pt}\setlength{\parsep}{0pt}}
\newcommand{\evidence}[2]{%
  \begin{minipage}{\linewidth}
  {\bfseries\textcolor{Magenta}{#1}}\\[0.25em]
  \centering\includegraphics[width=0.96\linewidth]{#2}
  \end{minipage}}

\begin{document}
\begin{frame}[t]

% ---------- Header ----------
\begin{beamercolorbox}[wd=\textwidth,sep=0.7cm,center,rounded=true]{block title}
{\fontsize{38}{43}\selectfont\bfseries From-Scratch Implementation and Quality Improvement of $\sinh(x)$}\\[0.25em]
{\fontsize{19}{23}\selectfont SOEN 6011 -- Software Engineering Processes \quad | \quad Delivery 3 (D3) \quad | \quad Function F3}\\[0.15em]
{\fontsize{17}{20}\selectfont Arvind Lakshmanan (40310757) \quad | \quad Concordia University \quad | \quad Summer 2026}
\end{beamercolorbox}
\vspace{0.45cm}

\begin{columns}[t,totalwidth=\textwidth]
% ==================== COLUMN 1 ====================
\begin{column}{0.318\textwidth}

\begin{block}{Project Evolution}
\small
\begin{center}
\begin{tikzpicture}[node distance=1.2cm, every node/.style={font=\small}]
\node[draw=Magenta,rounded corners=7pt,fill=LightPink,text width=7.0cm,inner sep=8pt] (d1) {
\textbf{D1}\\Persona; requirements; two pseudocode algorithms; algorithm-selection mind map; initial Python TUI.};
\node[draw=Orange,rounded corners=7pt,fill=LightOrange,text width=7.0cm,inner sep=8pt,right=0.9cm of d1] (d2) {
\textbf{D2}\\From-scratch implementation; Tkinter GUI; exception handling; public version control/README; updated requirements.};
\node[draw=DeepPurple,rounded corners=7pt,fill=LightPink,text width=7.0cm,inner sep=8pt,right=0.9cm of d2] (d3) {
\textbf{D3}\\PEP-8 improvement; Flake8; pdb; Pylint; Semantic Versioning 1.1.0; UIDP/accessibility; PyUnit tests.};
\draw[-{Latex[length=4mm]},thick,Magenta] (d1) -- (d2);
\draw[-{Latex[length=4mm]},thick,Orange] (d2) -- (d3);
\end{tikzpicture}
\end{center}
\end{block}

\begin{block}{Methodology}
\small
\justifying
The calculator computes $\sinh(x)$ \textbf{from scratch} using a recurrence-based Maclaurin series. The production calculation does not call \texttt{math.sinh}, \texttt{math.exp}, factorial helpers, NumPy, or SciPy.

\vspace{0.45em}
\begin{beamercolorbox}[rounded=true,sep=0.35cm,wd=0.96\linewidth]{block body alerted}
\centering
{\bfseries Maclaurin series}\\[0.2em]
$\displaystyle \sinh(x)=x+\frac{x^3}{3!}+\frac{x^5}{5!}+\cdots$
\end{beamercolorbox}

\vspace{0.35em}
\begin{beamercolorbox}[rounded=true,sep=0.35cm,wd=0.96\linewidth]{block body alerted}
\centering
{\bfseries Recurrence relation}\\[0.2em]
$\displaystyle t_0=x,\qquad t_k=t_{k-1}\frac{x^2}{(2k)(2k+1)}$
\end{beamercolorbox}

\vspace{0.45em}
\begin{center}
\begin{tikzpicture}[node distance=0.45cm, every node/.style={font=\scriptsize,align=center}]
\node[draw=Magenta,rounded corners,fill=LightPink,minimum width=3.1cm,minimum height=1.7cm] (a) {Input $x$\\validate $[-20,20]$};
\node[draw=Orange,rounded corners,fill=LightOrange,minimum width=3.1cm,minimum height=1.7cm,right=of a] (b) {Initialize\\$t=x,\;s=x$};
\node[draw=Orange,rounded corners,fill=LightOrange,minimum width=3.1cm,minimum height=1.7cm,right=of b] (c) {Iterate recurrence\\add to sum};
\node[draw=Green,rounded corners,fill=white,minimum width=3.1cm,minimum height=1.7cm,right=of c] (d) {Converged?\\relative tolerance};
\node[draw=DeepPurple,rounded corners,fill=LightPink,minimum width=3.1cm,minimum height=1.7cm,right=of d] (e) {Return\\$s\approx\sinh(x)$};
\foreach \x/\y in {a/b,b/c,c/d,d/e}{\draw[-{Latex[length=3mm]},thick] (\x)--(\y);}
\end{tikzpicture}
\end{center}

\vspace{0.35em}
\centering
\textbf{Stopping rule:} $|t_k|\le 10^{-15}(1+|s|)$ or maximum \textbf{200} terms.
\end{block}

\begin{block}{Correctness and Range}
\small
\begin{itemize}\tightitem
\item Supported input: one finite real number with $-20\le x\le20$.
\item Invalid, blank, NaN, infinity, and out-of-range inputs raise a helpful custom error.
\item The recurrence preserves the odd-function property $\sinh(-x)=-\sinh(x)$.
\item Version shown in source and GUI: \textbf{1.1.0}.
\end{itemize}
\end{block}

\begin{block}{GAI / CASTROFF Use}
\small
\textbf{GAI tool:} ChatGPT (OpenAI). Prompts were structured with CASTROFF:\\
\textbf{C}onstraints, \textbf{A}udience, \textbf{S}tructure, \textbf{T}one, \textbf{R}ole, \textbf{O}utput Format, \textbf{F}ocus, \textbf{F}unction.
\begin{itemize}\tightitem
\item \textbf{Problem 7:} compliance review/refinement for PEP-8, UIDP, accessibility, and evidence planning.
\item \textbf{Problem 8:} unit-test coverage review and PyUnit test design.
\item Real Flake8, Pylint, pdb, and unittest results were \textbf{verified from actual tool evidence}, not generated by GAI.
\end{itemize}
\textit{Full CASTROFF prompt examples and output explanations are provided in the accompanying GAI document.}
\end{block}

\end{column}
\hfill
% ==================== COLUMN 2 ====================
\begin{column}{0.334\textwidth}

\begin{alertblock}{Results \& Quality Evidence}
\small
\textbf{Flake8 style check}\par
\includegraphics[width=\linewidth]{Evidence/flake8.png}
\vspace{0.15em}
\tick\ Flake8 returned to the prompt with no listed style violations.

\vspace{0.45em}
\textbf{Pylint static analysis}\par
\includegraphics[width=\linewidth]{Evidence/pylint.png}
\vspace{0.15em}
\tick\ Pylint score: \textbf{9.86/10}, with two minor warnings shown in the evidence.

\vspace{0.45em}
\textbf{pdb debugger}\par
\includegraphics[width=\linewidth]{Evidence/pdb.png}
\vspace{0.15em}
\tick\ The debugger stepped through \texttt{calculate\_sinh\_from\_scratch} and inspected intermediate state.
\end{alertblock}

\begin{block}{GUI Design, UIDP \& Accessibility}
\scriptsize
\begin{columns}[t,totalwidth=\linewidth]
\begin{column}{0.43\linewidth}
\centering\includegraphics[width=\linewidth]{Evidence/gui.png}
\end{column}
\begin{column}{0.54\linewidth}
\begin{itemize}\tightitem
\item \textbf{Consistency:} uniform layout/controls.
\item \textbf{Visibility:} clear labels, results, status.
\item \textbf{Feedback:} textual success/error messages.
\item \textbf{Error prevention:} validates $-20\le x\le20$.
\item \textbf{Simplicity:} short calculation flow.
\item \textbf{User control:} Calculate, Clear, Exit, re-entry.
\item \textbf{Accessibility:} no color-only meaning.
\item \textbf{Keyboard:} Enter/Alt+C, Alt+L, Escape, Tab.
\end{itemize}
\end{column}
\end{columns}
\end{block}

\end{column}
\hfill
% ==================== COLUMN 3 ====================
\begin{column}{0.318\textwidth}

\begin{block}{Testing — PyUnit / unittest Evidence}
\small
\centering
\includegraphics[width=0.98\linewidth]{Evidence/tests.png}
\vspace{0.35em}
\begin{itemize}\tightitem
\item \textbf{Total tests: 16}
\item \textbf{Result: all tests passed (OK)}
\item Coverage includes helpers, parsing, valid/invalid inputs, NaN/infinity rejection, boundaries, known values, zero, and $\sinh(-x)=-\sinh(x)$.
\end{itemize}
\end{block}

\begin{alertblock}{UIDP Decision Mind Map}
\small
\centering
\begin{tikzpicture}[x=1cm,y=1cm, every node/.style={font=\scriptsize,align=center}]
\node[draw=DeepPurple,very thick,circle,minimum size=3.3cm,fill=LightPink] (center) at (0,0) {\bfseries UIDP\\for F3 GUI};

\node[draw=Green,rounded corners,fill=white,text width=4.0cm] (a1) at (-7.0,4.2) {Consistency};
\node[draw=Green,rounded corners,fill=white,text width=4.0cm] (a2) at (-7.0,3.0) {Visibility};
\node[draw=Green,rounded corners,fill=white,text width=4.0cm] (a3) at (-7.0,1.8) {Feedback};
\node[draw=Green,rounded corners,fill=white,text width=4.0cm] (a4) at (-7.0,0.6) {Error prevention};
\node[draw=Green,rounded corners,fill=white,text width=4.0cm] (a5) at (-7.0,-0.6) {Simplicity};
\node[draw=Green,rounded corners,fill=white,text width=4.0cm] (a6) at (-7.0,-1.8) {User control};
\node[draw=Green,rounded corners,fill=white,text width=4.0cm] (a7) at (-7.0,-3.0) {Accessibility};
\node[draw=Green,rounded corners,fill=white,text width=4.0cm] (a8) at (-7.0,-4.2) {Keyboard-friendly interaction};

\node[draw=Red,rounded corners,fill=white,text width=4.8cm] (n1) at (7.2,2.7) {Multi-user collaboration};
\node[draw=Red,rounded corners,fill=white,text width=4.8cm] (n2) at (7.2,0.9) {Complex customization};
\node[draw=Red,rounded corners,fill=white,text width=4.8cm] (n3) at (7.2,-0.9) {Advanced personalization};
\node[draw=Red,rounded corners,fill=white,text width=4.8cm] (n4) at (7.2,-2.7) {Gesture-based input};

\foreach \n in {a1,a2,a3,a4,a5,a6,a7,a8}{\draw[Green,thick] (\n.east)--(center.west);}
\foreach \n in {n1,n2,n3,n4}{\draw[Red,thick] (center.east)--(\n.west);}
\node[text=Green,font=\small\bfseries] at (-6.9,5.25) {Applies};
\node[text=Red,font=\small\bfseries] at (7.2,4.1) {Does not strongly apply};
\end{tikzpicture}

\vspace{0.4em}
\begin{beamercolorbox}[rounded=true,sep=0.3cm,wd=0.94\linewidth]{block body alerted}
\centering Chosen because the application is a \textbf{single-user calculator with a simple task flow}; principles were selected according to their relevance to this interface.
\end{beamercolorbox}
\end{alertblock}

\begin{block}{D3 Completion Summary}
\small
\begin{itemize}\tightitem
\item \tick\ PEP-8-oriented implementation and Flake8 evidence.
\item \tick\ pdb debugging evidence.
\item \tick\ Pylint static analysis evidence.
\item \tick\ Semantic Versioning 1.1.0.
\item \tick\ UIDP decision mind map and accessibility improvements.
\item \tick\ 16 PyUnit/unittest test cases, all passing in the submitted evidence.
\item \tick\ CASTROFF-based GAI documentation for Problems 7 and 8.
\end{itemize}
\end{block}

\end{column}
\end{columns}

\vfill
\begin{center}
\scriptsize
\textbf{References:} SOEN 6011 Project Description, Summer 2026; PEP 8; Python \texttt{unittest}; CASTROFF framework (Constraints, Audience, Structure, Tone, Role, Output Format, Focus, Function).
\end{center}

\end{frame}
\end{document}
